\section{从 3-Gram 到 GPT-3：语言模型的质量阶梯}

\term{language model}（语言模型）为 token 序列分配概率。给定 $x_1,\ldots,x_T$，自回归分解为
\[
p(x_1,\ldots,x_T)=\prod_{t=1}^{T}p(x_t\mid x_{<t}).
\]
模型差异主要在条件上下文如何表示、可使用多长历史、参数与数据如何扩展。讲者先展示 sample，而不是先展示 benchmark，是为了让读者直观看见概率模型的失败模式。

\subsection{3-Gram：局部统计能生成“相关词”，却难形成全局语义}

\term{n-gram} 假设当前 token 只依赖前 $n-1$ 个 token。3-gram 写为 $p(x_t\mid x_{t-2},x_{t-1})$。它可以从计数平滑得到概率，训练与解释简单，但上下文窗口固定，罕见组合稀疏，无法维护长距离实体、事实和篇章结构。

\lecturefigure{slide-01-3gram.jpg}{3-gram sample 已有局部相关词，却缺少稳定句法和全局语义。}{官方视频 00:08；Shannon 1951 historical framing。}

\begin{knowledgebox}{读图：不要只问句子“像不像英语”}
Sample 中金额、年份和商业词汇局部共现，说明短窗口统计捕获了词邻接；数字关系、实体一致性与句子目的却崩溃。评价时可分成 lexical coherence、syntactic coherence、semantic consistency、factual consistency 四层。n-gram 通常在前两层的短范围有信号，后两层缺乏状态。
\end{knowledgebox}

\begin{knowledgebox}{Perplexity 的定义与限制}
\term{perplexity}（困惑度）是平均 cross-entropy 的指数：
\[
\operatorname{PPL}=\exp\!\left(-\frac{1}{T}\sum_{t=1}^{T}\log p(x_t\mid x_{<t})\right).
\]
它可直觉理解为每步“有效候选数”，越低通常表示对该测试分布预测更好。但不同 tokenizer、数据清洗与 domain 的 PPL 不宜直接横比；更低 PPL 也不自动保证事实性、可控性或任务成功。
\end{knowledgebox}

\subsection{RNN：用隐藏状态替代固定窗口}

本节从固定窗口转向递归状态。RNN 更新 $h_t=f(x_t,h_{t-1})$，理论上能让任意历史影响当前预测。它消除了固定 n-gram 窗口，却把全部历史压入单个状态，并使训练沿时间串行。长序列还面临梯度消失/爆炸和信息覆盖。

\lecturefigure{slide-02-rnn.jpg}{RNN sample 比 3-gram 更连贯，但实体与长程结构仍容易漂移。}{官方视频 00:27；Sutskever et al. 2011。}

\begin{knowledgebox}{读图：RNN 改善的是状态连续性}
RNN 能维持更长的局部主题和字符/词模式，因此 sample 不再只像词袋。问题是每一步都通过同一个 recurrent bottleneck，旧信息要经过许多非线性变换才能到达未来。输出中出现“句子像语言、论证不像论证”的现象，说明流畅性与全局计划仍是不同能力。
\end{knowledgebox}

\teachervoice{讲者用真实 sample 逐级展示模型，而不是声称“新架构全面优于旧架构”。课堂提示：生成模型的进步常先表现为错误变得更难发现，因此 evaluation 也必须同步升级。}

\subsection{Big LSTM：门控和规模提高记忆，但没有消除串行路径}

\term{LSTM}（Long Short-Term Memory）用 input、forget、output gates 控制 cell state，缓解普通 RNN 的优化问题。更大 LSTM、更多数据和 better regularization 曾显著推进 language modeling，但训练依旧按时间步递归，长距离信息仍要穿过 $O(T)$ 路径。

\lecturefigure{slide-03-big-lstm.jpg}{大规模 LSTM 生成更像自然文本，但仍暴露重复、漂移和结构错误。}{官方视频 01:12；Jozefowicz et al. 2016。}

\begin{knowledgebox}{读图：scale 与 architecture 同时改变 sample}
Slide 将 Big LSTM 放在 RNN 与 Transformer 之间，提醒我们不要把所有进步归因于一个组件。门控改善记忆，模型/数据规模降低局部不确定性，训练 recipe 提高稳定性；但 recurrent dependency 仍限制硬件并行。公平比较需要控制参数量、token 数和 compute，而不是只选最好看的样本。
\end{knowledgebox}

\subsection{Transformer：并行训练与更短信息路径}

Transformer 用 self-attention 让 token 直接交互，矩阵运算更适合 accelerator，并允许扩展更大 batch、模型和数据。它没有自动解决事实、规划和目标问题，但改变了可行的 scaling regime。

\lecturefigure{slide-04-transformer.jpg}{Transformer sample 的提升来自架构、规模、数据与训练共同作用。}{官方视频 01:52。}

\begin{knowledgebox}{读图：Model completion 是条件分布，不是知识库查询}
Slide 给 prompt 和 completion，输出可以延续格式与主题，却仍可能编造细节。Language model 预测“在训练分布中什么续写可能”，不执行默认 fact-check。随着 sample 更流畅，人类更容易把语言质量误当事实质量，因此 later evaluation 需要任务 grounding、retrieval 或 external checks。
\end{knowledgebox}

\subsection{GPT-2：规模化 Transformer 与长文本样本}

本节转向 decoder-only scaling。GPT-2 将 Transformer 扩展到更大数据与参数，强调 zero-shot task behavior。课堂先展示模型与数据规模，再展示新闻样本，意图是让读者观察跨段 coherence、风格模仿和 failure。

\lecturefigure{slide-05-gpt2-model.jpg}{GPT-2 把 decoder-only Transformer 扩展为通用文本生成模型。}{官方视频 02:33；Radford et al. 2019。}

\begin{knowledgebox}{读图：参数量只是系统的一列}
模型规模扩大通常伴随更多训练 token、不同 batch/optimizer、context 和 compute。只说“参数更大所以能力出现”会忽略 data diversity 与 optimization。读这张 slide 时应把 architecture、parameters、data、objective 与 evaluation 作为共同 recipe。
\end{knowledgebox}

\lecturefigure{slide-06-gpt2-example-1.jpg}{GPT-2 新闻样本展示篇章风格与实体延续，同时保留可检查的虚构风险。}{官方视频 03:05。}

\begin{knowledgebox}{读图：长样本要按 claim 审核}
先看句法和段落连接，再标出人物、地点、数字与因果 claim；最后检查这些 claim 是否可由 prompt 或外部证据支持。模型可以保持新闻文体，却在细节处“自信地补齐”。样本评估若只问读起来顺不顺，会系统性漏掉 hallucination。
\end{knowledgebox}

\lecturefigure{slide-07-gpt2-example-2.jpg}{第二个 GPT-2 样本说明不同 prompt 会暴露不同类型的 coherence failure。}{官方视频 03:38。}

\begin{knowledgebox}{读图：不要用单一样本代表模型分布}
Sampling seed、temperature、top-p 与 prompt 都会改变输出。一个漂亮样本不能代表平均性能，一个糟糕样本也不能证明模型无用。需要固定 decoding protocol、报告样本数量，并结合 likelihood、human rating 与 downstream success。
\end{knowledgebox}

\subsection{GPT-3：人类能否识别生成新闻}

当文本质量提高，evaluation 从“看起来很差”转向受控人类实验。GPT-3 论文让参与者判断新闻是否由模型生成；准确率随模型规模变化。这个实验测的是特定 prompt、长度、采样和参与者条件下的 detection，不等同真实性或通用智能。

\lecturefigure{slide-08-gpt3-human-detection.jpg}{GPT-3 用人类识别任务把 sample realism 转成可量化实验。}{官方视频 05:12；Brown et al. 2020。}

\begin{knowledgebox}{读图：先定义实验单位}
每个 item 是模型生成或人类撰写的新闻片段；参与者输出二分类判断。结果受主题熟悉度、文本长度、生成温度与 instruction 影响。接近 chance 只能说明该设置下难以区分来源，不说明内容真实、无害或可用于新闻生产。
\end{knowledgebox}

\lecturefigure{slide-09-gpt3-detection-curve.jpg}{更大模型使检测更困难，但曲线仍需与实验设置一起解释。}{官方视频 08:10；Brown et al. 2020。}

\begin{knowledgebox}{读图：横轴是规模，纵轴不是“智能”}
曲线把 model size 与 human detection accuracy 联系起来。关键趋势是规模增大后生成文体更接近人类分布；它不测事实性、来源可追溯或长期一致性。误差条、baseline 与 sample policy 决定结论强度。不能把单一曲线外推到所有语言和所有生成任务。
\end{knowledgebox}

\subsection{本章小结}

从 n-gram 到 GPT-3，sample 逐步变得连贯，评价也从肉眼失败升级到受控人类实验。Scale 改变分布质量，却让真实性与来源识别问题更加重要。

